\chapter{End Your Day in a Hot Spring}

% ================= PERCAKAPAN 1-2 =================
\begin{convobox}{1}
    \noindent
    \jpkana{こんばんは。}{Konbanwa.}

    \vspace{1.5em}
    \noindent{\Large \color{articolor} \textbf{Arti:} Selamat malam.}
\end{convobox}

\begin{convobox}{2}
    \noindent
    \jpword{きょう}{今日}{Kyou} \jpkana{も}{mo} 
    \jpkana{お}{o}\jpword{つか}{疲}{tsuka}\jpkana{れ}{re}\jpword{さま}{様}{sama} \jpkana{でした。}{deshita.}

    \vspace{1.5em}
    \noindent{\Large \color{articolor} \textbf{Arti:} Hari ini juga, terima kasih atas kerja kerasmu.}
\end{convobox}

\begin{lessonbox}{Catatan Belajar: Selamat Malam!}
    Adik-adik, kalau hari sudah gelap, kita sapa teman dengan \textbf{こんばんは (Konbanwa)} yang artinya "Selamat malam". Terus, perhatikan kata \textit{Otsukaresama deshita}. Di bab sebelumnya pakai \textit{desu}, sekarang pakai \textit{deshita} karena harinya sudah hampir berakhir!
\end{lessonbox}

% ================= PERCAKAPAN 3-4 =================
\begin{convobox}{3}
    \noindent
    \jpkana{ゆっくり}{Yukkuri} 
    \jpword{おんせん}{温泉}{onsen} \jpkana{に}{ni} 
    \jpword{つか}{浸}{tsuka}\jpkana{りながら、}{ri nagara,}

    \vspace{1.5em}
    \noindent{\Large \color{articolor} \textbf{Arti:} Sambil berendam santai di pemandian air panas (onsen),}
\end{convobox}

\begin{convobox}{4}
    \noindent
    \jpword{いちにち}{一日}{Ichi-nichi} \jpkana{を}{o} 
    \jpword{ふ}{振}{fu}\jpkana{り}{ri}\jpword{かえ}{返}{kae}\jpkana{りましょう。}{rimashou.}

    \vspace{1.5em}
    \noindent{\Large \color{articolor} \textbf{Arti:} mari kita ingat-ingat (renungkan) apa saja yang terjadi hari ini.}
\end{convobox}

\begin{lessonbox}{Catatan Belajar: Melakukan 2 Hal Sekaligus!}
    Coba lihat kata \textbf{〜ながら (〜nagara)}. Ini artinya "Sambil". Jadi kalau kamu mandi sambil nyanyi, atau makan sambil nonton TV, kamu bisa pakai \textit{nagara} di tengah kalimatnya. Seru kan?
\end{lessonbox}

% ================= PERCAKAPAN 5-6 =================
\begin{convobox}{5}
    \noindent
    \jpword{きょう}{今日}{Kyou} \jpkana{の}{no} 
    \jpword{てんき}{天気}{tenki} \jpkana{は}{wa} 
    \jpkana{どうでしたか？}{dou deshita ka?}

    \vspace{1.5em}
    \noindent{\Large \color{articolor} \textbf{Arti:} Bagaimana cuaca hari ini?}
\end{convobox}

\begin{convobox}{6}
    \noindent
    \jpword{いちにちじゅう}{一日中}{Ichi-nichi-juu} 
    \jpword{あめ}{雨}{ame} \jpkana{でした。}{deshita.}

    \vspace{1.5em}
    \noindent{\Large \color{articolor} \textbf{Arti:} Sepanjang hari hujan.}
\end{convobox}

\begin{lessonbox}{Catatan Belajar: Sepanjang Hari}
    Kalau kamu melakukan sesuatu dari pagi sampai malam, kamu bisa bilang \textbf{一日中 (Ichi-nichi-juu)}. Di cerita ini hujannya \textit{ichi-nichi-juu} alias nggak berhenti seharian!
\end{lessonbox}

% ================= PERCAKAPAN 7-8 =================
\begin{convobox}{7}
    \noindent
    \jpkana{そうですか。}{Sou desu ka.}

    \vspace{1.5em}
    \noindent{\Large \color{articolor} \textbf{Arti:} Oh, begitu.}
\end{convobox}

\begin{convobox}{8}
    \noindent
    \jpword{きょう}{今日}{Kyou} \jpkana{は}{wa} 
    \jpword{いそが}{忙}{isoga}\jpkana{しかったですか？}{shikatta desu ka?}

    \vspace{1.5em}
    \noindent{\Large \color{articolor} \textbf{Arti:} Apakah hari ini sibuk?}
\end{convobox}

\begin{lessonbox}{Catatan Belajar: Tanya Kabar Masa Lalu}
    Ingat kan kalau kata sifat bisa berubah jadi masa lalu? Kata "Sibuk" adalah \textit{Isogashii}. Karena bertanyanya malam hari untuk kejadian siang tadi, kata \textit{Isogashii} berubah bentuk menjadi \textbf{Isogashikatta}!
\end{lessonbox}

% ================= PERCAKAPAN 9-10 =================
\begin{convobox}{9}
    \noindent
    \jpkana{はい、}{Hai,} 
    \jpword{すこ}{少}{suko}\jpkana{し}{shi} 
    \jpword{いそが}{忙}{isoga}\jpkana{しかったです。}{shikatta desu.}

    \vspace{1.5em}
    \noindent{\Large \color{articolor} \textbf{Arti:} Ya, sedikit sibuk.}
\end{convobox}

\begin{convobox}{10}
    \noindent
    \jpword{きょう}{今日}{Kyou} \jpkana{は、}{wa,} 
    \jpword{なに}{何}{nani} \jpkana{を}{o} 
    \jpkana{しましたか？}{shimashita ka?}

    \vspace{1.5em}
    \noindent{\Large \color{articolor} \textbf{Arti:} Hari ini, (kamu) melakukan apa?}
\end{convobox}

\begin{lessonbox}{Catatan Belajar: "Melakukan Apa?"}
    Ini adalah pertanyaan penting kalau kamu mau kepo sama hari temanmu! \textbf{何をしましたか？ (Nani o shimashita ka?)}. Kamu bisa jawab dengan aktivitas yang kamu lakukan hari itu.
\end{lessonbox}

% ================= PERCAKAPAN 11-12 =================
\begin{convobox}{11}
    \noindent
    \jpkana{えーっと、}{Eetto,} 
    \jpword{せんたく}{洗濯}{sentaku} \jpkana{を}{o} \jpkana{して、}{shite,} 
    \jpword{りょうり}{料理}{ryouri} \jpkana{を}{o} \jpkana{して…}{shite...}

    \vspace{1.5em}
    \noindent{\Large \color{articolor} \textbf{Arti:} Emm, mencuci baju, lalu memasak...}
\end{convobox}

\begin{convobox}{12}
    \noindent
    \jpkana{その}{Sono} \jpword{あと}{後}{ato} 
    \jpkana{カフェ}{kafe} \jpkana{に}{ni} 
    \jpword{い}{行}{i}\jpkana{って}{tte} 
    \jpword{ほん}{本}{hon} \jpkana{を}{o} 
    \jpword{よ}{読}{yo}\jpkana{みました。}{mimashita.}

    \vspace{1.5em}
    \noindent{\Large \color{articolor} \textbf{Arti:} setelah itu pergi ke kafe dan membaca buku.}
\end{convobox}

\begin{lessonbox}{Catatan Belajar: Menyambung Kegiatan}
    Kalau kamu melakukan banyak hal berturut-turut, kamu bisa ubah akhirannya jadi \textbf{〜て (〜te)}. Seperti di atas: \textit{Sentaku o shi\textbf{te}}, \textit{Ryouri o shi\textbf{te}}. Artinya "Melakukan A, lalu melakukan B".
\end{lessonbox}

% ================= PERCAKAPAN 13-14 =================
\begin{convobox}{13}
    \noindent
    \jpword{わたし}{私}{Watashi} \jpkana{は、}{wa,} 
    \jpword{いちにちじゅう}{一日中}{ichi-nichi-juu} 
    \jpword{しごと}{仕事}{shigoto} \jpkana{を}{o} \jpkana{して}{shite}

    \vspace{1.5em}
    \noindent{\Large \color{articolor} \textbf{Arti:} Kalau aku, seharian bekerja,}
\end{convobox}

\begin{convobox}{14}
    \noindent
    \jpkana{ちょっと}{chotto} 
    \jpword{つか}{疲}{tsuka}\jpkana{れたので、}{reta node,} 
    \jpword{おんせん}{温泉}{onsen} \jpkana{に}{ni} 
    \jpword{き}{来}{ki}\jpkana{ました。}{mashita.}

    \vspace{1.5em}
    \noindent{\Large \color{articolor} \textbf{Arti:} karena sedikit lelah, jadi aku datang ke onsen.}
\end{convobox}

\begin{lessonbox}{Catatan Belajar: "Karena..." (Sekali Lagi!)}
    Di Bab 2 kita sudah bahas ini loh. \textbf{〜ので (〜node)} artinya "Karena". Di kalimat ini teman kita bilang \textit{Tsukareta node} yang artinya "Karena aku lelah". Makanya dia pergi berendam deh!
\end{lessonbox}

% ================= PERCAKAPAN 15-16 =================
\begin{convobox}{15}
    \noindent
    \jpword{おんせん}{温泉}{Onsen} \jpkana{に}{ni} 
    \jpword{はい}{入}{hai}\jpkana{ると、}{ru to,} 
    \jpword{つか}{疲}{tsuka}\jpkana{れが}{re ga} 
    \jpkana{よく}{yoku} 
    \jpword{と}{取}{to}\jpkana{れるんですよ。}{reru n desu yo.}

    \vspace{1.5em}
    \noindent{\Large \color{articolor} \textbf{Arti:} Kalau masuk ke onsen, rasa lelahnya bisa hilang lho.}
\end{convobox}

\begin{convobox}{16}
    \noindent
    \jpword{つか}{疲}{Tsuka}\jpkana{れた}{reta} 
    \jpword{とき}{時}{toki} \jpkana{や、}{ya,} 
    \jpkana{リラックス}{rirakkusu} \jpkana{したい}{shitai} 
    \jpword{とき}{時}{toki}\jpkana{、}{,}

    \vspace{1.5em}
    \noindent{\Large \color{articolor} \textbf{Arti:} Saat merasa lelah, atau saat ingin bersantai,}
\end{convobox}

\begin{lessonbox}{Catatan Belajar: Kata "Saat / Waktu"}
    Dalam bahasa Jepang, kalau kita mau bilang "Saat... / Ketika...", kita memakai kata \textbf{時 (Toki)}. Misalnya, \textit{Tsukareta toki} (Saat lelah). Mudah, kan?
\end{lessonbox}

% ================= PERCAKAPAN 17-18 =================
\begin{convobox}{17}
    \noindent
    \jpkana{よく}{Yoku} 
    \jpword{なに}{何}{nani} \jpkana{を}{o} 
    \jpkana{しますか？}{shimasu ka?}

    \vspace{1.5em}
    \noindent{\Large \color{articolor} \textbf{Arti:} biasanya kamu sering melakukan apa?}
\end{convobox}

\begin{convobox}{18}
    \noindent
    \jpkana{ハーブティー}{Haabutii} \jpkana{を}{o} 
    \jpword{の}{飲}{no}\jpkana{んだり、}{ndari,}

    \vspace{1.5em}
    \noindent{\Large \color{articolor} \textbf{Arti:} Meminum teh herbal, atau}
\end{convobox}

\begin{lessonbox}{Catatan Belajar: Grammar Favorit Kita!}
    Yey, kita ketemu lagi sama grammar \textbf{〜たり、〜たり (〜tari, 〜tari)}! Ini untuk menyebutkan hal acak yang kita lakukan. Salah satunya adalah \textit{Nondari} (Minum), berasal dari kata \textit{Nomimasu}.
\end{lessonbox}

% ================= PERCAKAPAN 19-20 =================
\begin{convobox}{19}
    \noindent
    \jpword{さんぽ}{散歩}{Sanpo} \jpkana{に}{ni} 
    \jpword{い}{行}{i}\jpkana{ったりします。}{ttari shimasu.}

    \vspace{1.5em}
    \noindent{\Large \color{articolor} \textbf{Arti:} pergi jalan-jalan santai.}
\end{convobox}

\begin{convobox}{20}
    \noindent
    \jpkana{いいですね。}{Ii desu ne.}

    \vspace{1.5em}
    \noindent{\Large \color{articolor} \textbf{Arti:} Wah bagus ya.}
\end{convobox}

\begin{lessonbox}{Catatan Belajar: "Wah, Bagus Ya!"}
    Kalau ada temanmu bercerita hal yang menyenangkan, kamu bisa memujinya dengan senyuman dan bilang \textbf{いいですね (Ii desu ne!)} yang berarti "Bagus/Hebat ya!".
\end{lessonbox}

% ================= PERCAKAPAN 21-22 =================
\begin{convobox}{21}
    \noindent
    \jpkana{ところで、}{Tokoro de,} 
    \jpword{しゅみ}{趣味}{shumi} \jpkana{は}{wa} 
    \jpword{なに}{何}{nani}\jpkana{か}{ka} 
    \jpkana{ありますか？}{arimasu ka?}

    \vspace{1.5em}
    \noindent{\Large \color{articolor} \textbf{Arti:} Ngomong-ngomong, apakah kamu punya suatu hobi?}
\end{convobox}

\begin{convobox}{22}
    \noindent
    \jpkana{ギター}{Gitaa} \jpkana{を}{o} 
    \jpword{ひ}{弾}{hi}\jpkana{くことです。}{ku koto desu.}

    \vspace{1.5em}
    \noindent{\Large \color{articolor} \textbf{Arti:} (Hobiku) adalah bermain gitar.}
\end{convobox}

\begin{lessonbox}{Catatan Belajar: Mengganti Topik \& Hobi}
    Kalau mau pindah ke topik obrolan baru, gunakan \textbf{ところで (Tokoro de)} yang artinya "Ngomong-ngomong". Untuk menjawab hobi, kita gunakan kata \textbf{こと (koto)} setelah kegiatannya supaya jadi kata benda. Contohnya \textit{Hiku koto} (Bermain alat musik).
\end{lessonbox}

% ================= PERCAKAPAN 23-24 =================
\begin{convobox}{23}
    \noindent
    \jpkana{そうなんですね。}{Sou nan desu ne.}

    \vspace{1.5em}
    \noindent{\Large \color{articolor} \textbf{Arti:} Oh begitu ya.}
\end{convobox}

\begin{convobox}{24}
    \noindent
    \jpkana{いつ}{Itsu} 
    \jpkana{その}{sono} 
    \jpword{しゅみ}{趣味}{shumi} \jpkana{を}{o} 
    \jpword{はじ}{始}{haji}\jpkana{めたんですか？}{meta n desu ka?}

    \vspace{1.5em}
    \noindent{\Large \color{articolor} \textbf{Arti:} Kapan kamu mulai hobi tersebut?}
\end{convobox}

\begin{lessonbox}{Catatan Belajar: "Sou nan desu ne!"}
    Kata ini mirip seperti \textit{Sou desu ka}, tapi mendengarnya terasa lebih hangat dan antusias! \textbf{そうなんですね (Sou nan desu ne)} biasa dipakai saat kita mendengarkan cerita teman dengan penuh perhatian.
\end{lessonbox}

% ================= PERCAKAPAN 25-26 =================
\begin{convobox}{25}
    \noindent
    \jpword{にねんまえ}{2年前}{Ni-nen mae} 
    \jpkana{ぐらいです。}{gurai desu.}

    \vspace{1.5em}
    \noindent{\Large \color{articolor} \textbf{Arti:} Kira-kira 2 tahun yang lalu.}
\end{convobox}

\begin{convobox}{26}
    \noindent
    \jpword{きょう}{今日}{Kyou} \jpkana{は、}{wa,} 
    \jpkana{それを}{sore o} 
    \jpkana{しましたか？}{shimashita ka?}

    \vspace{1.5em}
    \noindent{\Large \color{articolor} \textbf{Arti:} Hari ini, apakah kamu melakukannya?}
\end{convobox}

\begin{lessonbox}{Catatan Belajar: "Sore" artinya "Itu"}
    Daripada mengulang kata "bermain gitar" berulang kali, bahasa Jepang punya kata tunjuk \textbf{それ (Sore)} yang artinya "Itu". Jadi \textit{Sore o shimashita ka?} berarti "Apakah kamu melakukan (hobi) itu?".
\end{lessonbox}

% ================= PERCAKAPAN 27-28 =================
\begin{convobox}{27}
    \noindent
    \jpkana{はい、}{Hai,} 
    \jpword{いちじかん}{1時間}{ichi-jikan} \jpkana{ぐらい}{gurai} 
    \jpkana{ギター}{gitaa} \jpkana{を}{o} 
    \jpword{ひ}{弾}{hi}\jpkana{きました。}{kimashita.}

    \vspace{1.5em}
    \noindent{\Large \color{articolor} \textbf{Arti:} Ya, aku bermain gitar kira-kira 1 jam.}
\end{convobox}

\begin{convobox}{28}
    \noindent
    \jpword{きょう}{今日}{Kyou} \jpkana{は、}{wa,} 
    \jpword{なんじ}{何時}{nan-ji} \jpkana{に}{ni} 
    \jpword{ね}{寝}{ne}\jpkana{るつもりですか？}{ru tsumori desu ka?}

    \vspace{1.5em}
    \noindent{\Large \color{articolor} \textbf{Arti:} Hari ini, kamu berencana mau tidur jam berapa?}
\end{convobox}

\begin{lessonbox}{Catatan Belajar: Rencana / Niat}
    Kita bisa menyatakan niat kita melakukan sesuatu pakai kata \textbf{〜つもり (〜tsumori)}. Kalau kamu ditanya \textit{Neru tsumori desu ka?}, artinya temanmu nanya "Kamu ada niatan/rencana mau tidur jam berapa?".
\end{lessonbox}

% ================= PERCAKAPAN 29-30 =================
\begin{convobox}{29}
    \noindent
    \jpword{じゅうじ}{10時}{Juu-ji} \jpkana{に}{ni} 
    \jpword{ね}{寝}{ne}\jpkana{るつもりです。}{ru tsumori desu.}

    \vspace{1.5em}
    \noindent{\Large \color{articolor} \textbf{Arti:} Rencananya mau tidur jam 10.}
\end{convobox}

\begin{convobox}{30}
    \noindent
    \jpword{わたし}{私}{Watashi} \jpkana{は}{wa} 
    \jpword{ふだん}{普段}{fudan} 
    \jpword{じゅういちじ}{11時}{juu-ichi-ji} \jpkana{に}{ni} 
    \jpword{ね}{寝}{ne}\jpkana{ますが、}{masu ga,}

    \vspace{1.5em}
    \noindent{\Large \color{articolor} \textbf{Arti:} Kalau aku biasanya tidur jam 11, tapi}
\end{convobox}

\begin{lessonbox}{Catatan Belajar: "Biasanya"}
    Adik-adik, kalau mau cerita kebiasaan rutinitasmu, pakai kata \textbf{普段 (Fudan)} yang artinya "Biasanya". Misalnya: \textit{Fudan wa 9-ji ni nemasu} (Biasanya aku tidur jam 9).
\end{lessonbox}

% ================= PERCAKAPAN 31-33 =================
\begin{convobox}{31}
    \noindent
    \jpword{きょう}{今日}{Kyou} \jpkana{は}{wa} 
    \jpword{すこ}{少}{suko}\jpkana{し}{shi} 
    \jpword{はや}{早}{haya}\jpkana{めに}{me ni} 
    \jpword{ね}{寝}{ne}\jpkana{ようと}{you to} 
    \jpword{おも}{思}{omo}\jpkana{います。}{imasu.}

    \vspace{1.5em}
    \noindent{\Large \color{articolor} \textbf{Arti:} hari ini sepertinya aku ingin/pikir mau tidur lebih awal.}
\end{convobox}

\begin{convobox}{32}
    \noindent
    \jpkana{お}{o}\jpword{はなし}{話}{hanashi} \jpkana{できて、}{dekite,} 
    \jpword{たの}{楽}{tano}\jpkana{しかったです。}{shikatta desu.}

    \vspace{1.5em}
    \noindent{\Large \color{articolor} \textbf{Arti:} Bisa mengobrol (denganmu), rasanya menyenangkan.}
\end{convobox}

\begin{convobox}{33}
    \noindent
    \jpkana{それでは、}{Sore dewa,} 
    \jpkana{また}{mata} 
    \jpword{あした}{明日}{ashita}\jpkana{。}{.}

    \vspace{1.5em}
    \noindent{\Large \color{articolor} \textbf{Arti:} Kalau begitu, sampai jumpa besok.}
\end{convobox}

\begin{lessonbox}{Catatan Belajar: Berpisah di Malam Hari}
    Saat mau tidur atau berpamitan di ujung hari, kita bisa pakai \textbf{また明日 (Mata ashita)} yang artinya "Sampai jumpa besok!". Kamu juga bisa tidur lebih awal atau \textit{Hayame ni neru} biar besok pagi bangun dengan segar!
\end{lessonbox}

% --- SUMMARY BOX ---
\vspace{2em}
\begin{tcolorbox}[
    colback=blue!5!white,
    colframe=blue!80!black,
    boxrule=3pt,
    arc=5mm,
    title={\huge\bfseries \sffamily 🌟 Rangkuman Bab 7 🌟},
    fonttitle=\color{white},
    attach boxed title to top center={yshift=-4mm},
    boxed title style={colback=blue!80!black, rounded corners},
    left=5mm, right=5mm, top=7mm, bottom=5mm
]
\Large \textbf{Luar Biasa! Kita belajar banyak sekali malam ini!} Ini dia ringkasannya:
\begin{itemize}
    \item \textbf{Menyambung Kalimat:} Gunakan \textbf{〜て (〜te)} untuk kegiatan berturut-turut, dan \textbf{〜ながら (〜nagara)} untuk kegiatan yang dilakukan "Sambil".
    \item \textbf{Kata Keterangan:} \textbf{Ichi-nichi-juu} (Seharian), \textbf{Fudan} (Biasanya), dan \textbf{Tokoro de} (Ngomong-ngomong).
    \item \textbf{Waktu / Saat:} Gunakan \textbf{〜時 (〜toki)}.
    \item \textbf{Menanyakan Niat/Rencana:} Pakai \textbf{〜つもり (〜tsumori)}.
    \item \textbf{Salam:} \textbf{Konbanwa} (Malam), dan \textbf{Mata ashita} (Sampai besok).
\end{itemize}
Karena kita sudah belajar dengan giat dari tadi, yuk kita \textit{yukkuri} beristirahat dan jangan lupa \textit{hayame ni neyou} (tidur lebih cepat) ya! Mata ashita!
\end{tcolorbox}

\newpage